\section{Vulnerability Detection Methodologies}

This section defines the benchmark methodology, the detection logic of each paradigm and the evaluation criteria applied to the resulting outputs.

\subsection{Methodology}

The dataset is built on the Minimal Viable Exploit (MVE) principle, where each OWASP category is represented by one or two pairs of contracts, each Vulnerable contract has a syntactically equivalent Fixed counterpart that mitigates the specific OWASP vector and changes nothing else. The dataset contains 30 Solidity contracts (15 Vulnerable + 15 Fixed) mapped to the OWASP Smart Contract Top 10. Pairing makes both true positives (on Vulnerable contracts) and false positives (on Fixed contracts) directly measurable on the same code surface.

\textbf{Detection criteria}. A finding is counted as a True Positive only when it identifies the specific OWASP vector instantiated by the MVE — for example, a reentrancy-eth finding on \texttt{MVE\_O8\_Vuln.sol} counts; a generic low-level-call warning on the same contract does not. Any vulnerability meeting the per-tool severity thresholds below reported in a Fixed contract is classified as a False Positive. Detection criteria were normalized across tools to filter out non-critical noise:
\begin{itemize}
  \item Static Analysis (Slither, Aderyn): High or Medium severity findings only.
  \item Symbolic Execution (Mythril): Any identified SWC (Smart Contract Weakness Classification) code.
  \item Fuzzing (Echidna, Medusa): Any failed invariant.
  \item LLM Review: Scores $\leq 4$ on Vulnerable contracts; scores $\geq 6$ on Fixed contracts.
\end{itemize}

\textbf{Tool selection}. Selection prioritized production-grade analyzers in active maintenance and broad industry use. Tools requiring incompatible environments were excluded — Manticore depends on pysha3 and has been unmaintained since 2022; SmartCheck requires Java 8 and breaks under modern JDKs (javax.xml.bind removal). Semgrep and Solhint are included as standard-tooling baselines: they represent the lowest-friction analyzers a Solidity developer encounters in a default CI pipeline and their detection rate measures what protection the standard toolchain offers without dedicated security tooling.

\textbf{Reproducibility}. All tool versions and configuration parameters are pinned in \texttt{config/versions.env} and per-tool configuration files. Table~\ref{tab:tool-versions} summarizes the values used for the benchmark runs.

\begin{table}[h]
\centering
\caption{Tool versions and key parameters.}
\label{tab:tool-versions}
\resizebox{\columnwidth}{!}{%
\begin{tabular}{lll}
\toprule
\textbf{Tool} & \textbf{Version} & \textbf{Key parameters} \\
\midrule
Slither          & 0.10.4       & default detectors, JSON output \\
Aderyn           & 0.6.8        & per-file scope (\texttt{-i}) \\
Semgrep          & 1.75.0       & \texttt{p/smart-contracts} ruleset \\
Solhint          & 5.0.3        & \texttt{solhint:recommended} \\
Mythril          & 0.24.8       & execution timeout 60s \\
Halmos           & 0.2.1        & default symbolic depth \\
Echidna          & 2.3.2        & \texttt{testLimit}$^{*}$\texttt{=50000}, \texttt{seqLen}$^{*}$\texttt{=100} \\
Medusa           & 1.5.1        & \texttt{testLimit=50000}, \texttt{workers}$^{*}$\texttt{=4} \\
LLaMA-3.3-70B    & via Groq     & \texttt{temperature=0.1}, \texttt{max\_tokens=2048} \\
Gemini 2.5 Flash & google-genai & \texttt{temperature=0.1} \\
Qwen3-32B        & via Groq     & \texttt{temperature=0.1}, \texttt{max\_tokens=4096} \\
\bottomrule
\end{tabular}%
}\\[3pt]
\parbox{\columnwidth}{\scriptsize$^{*}$\texttt{testLimit}: randomized transactions per run; \texttt{seqLen}: max transaction sequence length; \texttt{workers}: parallel fuzzing processes.}
\end{table}

\FloatBarrier

\textbf{Pipeline}. The execution pipeline is sequential: static analysis (Slither, Aderyn, Semgrep, Solhint), symbolic and formal analysis (Mythril, Halmos), property-based fuzzing (Echidna, Medusa) and LLM-assisted review (LLaMA, Gemini, Qwen). Each stage writes outputs to a deterministic results directory; a unified report is generated by \texttt{scripts/build\_report.py}.

\subsection{Detection paradigms}

Each paradigm has its own internal logic, false-positive profile and operational cost. The remainder of this section describes how each paradigm reaches a verdict and where its structural limits lie.

\subsubsection{Static analysis: source-level pattern detection}

Static analyzers read the Solidity source and flag suspicious patterns without executing the code. The paradigm splits cleanly into two subcategories that produce structurally different output and should not be aggregated when reporting recall.

Semantic analyzers (Slither, Aderyn) build a Solidity-specific intermediate representation of the contract and apply detectors that encode EVM-specific knowledge. Slither's reentrancy-eth detector, for instance, recognizes the call-before-state-write pattern that defines CWE-841. Aderyn's tx-origin detector recognizes authentication-by-origin as a phishing surface. The detectors know what they are looking for at the level of attack vector.

Pattern matchers and linters (Semgrep, Solhint) operate without domain knowledge. Semgrep matches generic syntactic templates from a community ruleset. Solhint enforces style conventions and a small set of surface-level security rules. They are not vulnerability-aware in the same sense — they detect what their author wrote a rule for, not what the underlying exploit pattern requires.

\subsubsection{Symbolic execution and formal verification: path-level reasoning}

Both Mythril and Halmos reason about contract execution paths using SMT solvers, but they answer different questions and require different inputs from the user.

Mythril runs autonomously on EVM bytecode and checks each execution path against a fixed catalogue of vulnerability templates from the SWC registry. When the solver finds an input that drives the contract along a path violating one of these templates, Mythril reports a finding. The user provides only the contract and Mythril provides the templates.

Halmos does the opposite: the user provides the property to prove, written as a Solidity function prefixed \texttt{check\_*}. Halmos then verifies that property over the entire symbolic input space. A successful execution does not simply mean that ``no counterexample was found within the search frontier'', it is a formal proof that no input exists for which the property fails. The cost is that the user must author each property, but the benefit is that a proof, when obtained, is conclusive.

\subsubsection{Property-based fuzzing: invariant testing through randomized execution}

Echidna and Medusa generate random sequences of transactions and check, after each one, whether user-defined invariants still hold. The invariants are functions prefixed \texttt{echidna\_} that return a boolean — for example, \texttt{echidna\_balance\_never\_decreases()}. If the fuzzer finds a transaction sequence that breaks the invariant, it reports a minimized counterexample.

The two fuzzers share the same invariant language but differ in implementation. Echidna is single-process and written in Haskell; Medusa parallelizes workers natively in Go. Both are included in the benchmark because cross-validation matters: when Echidna and Medusa report identical findings on a contract, that result is unlikely to be specific to one fuzzer's implementation.

\subsubsection{LLM-based analysis: prompt-based audit}

The three models receive the contract source through a single API call with a fixed audit prompt (Listing~\ref{lst:llm-prompt}). The prompt asks for a structured response: vulnerability identification, severity classification and a single overall security score on a 1–10 scale (1 = critically vulnerable, 10 = perfectly secure). Temperature is set to 0.1 across all three models to reduce response variability across runs.

The score is extracted from the structured response field. For each model, three properties characterize performance on the benchmark:

\begin{itemize}
  \item \textbf{Discrimination} ($\Delta$ = avg(Fixed) $-$ avg(Vuln)): whether the model assigns lower scores to Vulnerable contracts and higher scores to Fixed variants on the same code surface.
  \item \textbf{False-positive rate on Fixed}: whether the model recognizes mitigation patterns (SafeERC20, TWAP oracles, signature nonces, \texttt{\_disableInitializers}) rather than scoring on residual code complexity.
  \item \textbf{Output parseability}: whether the structured response format is preserved across all 30 contracts.
\end{itemize}

The empirical values are reported in Section~6.

\begin{lstlisting}[
  caption={Audit prompt sent to all three LLMs.},
  label={lst:llm-prompt},
  basicstyle=\footnotesize\ttfamily,
  frame=single,
  breaklines=true
]
You are an expert Solidity smart contract security auditor. Analyze the following smart contract and identify ALL security vulnerabilities

For each vulnerability found, report:
 1. Vulnerability type (e.g., reentrancy, integer overflow, access control, etc.)
 2. Severity: Critical / High / Medium / Low / Informational
 3. Affected function and line number(s)
 4. Brief description of the risk
 5. Recommended fix (1-2 sentences)

Also provide at the end:
 - Overall security score (1-10, where 10 = perfectly secure, 1 = critically vulnerable)
 - One-sentence summary of the most critical finding (or "No critical issues found")

Contract:
```solidity
{code}
```
\end{lstlisting}

\par\smallskip
\noindent\textbf{Implementation note:} The Qwen3-32B invocation prepends the literal prefix \texttt{/no\_think\textbackslash n\textbackslash n} to the prompt (see \texttt{scripts/llm\_analyze.py});

LLM evidence is treated as a semantic complement, not as execution-ground truth. The paradigm has structural limits — no direct execution of EVM paths, susceptibility to hallucinated findings, and dependence on training distribution and prompt framing — that prevent it from substituting for static, symbolic or fuzzing analysis.

Section 6 reports the empirical results of applying these paradigms to the 30-contract dataset.